\exercisegroup

* 1. Let E be a topological space and let Σ be one of the species of structures defined in Exercises 7 and 8 of §2. Take the morphisms to be those defined in the same exercises, and the α-mappings to be the continuous mappings of E into a Σ-set. Show that the universal mapping problem for E (relative to the preceding definitions) has in general no solution. *

* 2. Let E be a field, Σ the species of algebraically closed field structures. Take the morphisms to be homomorphisms, and the α-mappings to be homomorphisms of E into an algebraically closed field. Show that the algebraic closure \(F_{E}\) of E and the canonical injection of E into \(F_{E}\) satisfy \((AU_{I}^{\prime})\) , but that there is in general no solution to the universal mapping problem for E. *

\begin{enumerate}
\def\labelenumi{\arabic{enumi}.}
\setcounter{enumi}{2}
\tightlist
\item
  Let \(\Sigma\) be a species of structures, let \((\mathbf{A}_i)_{i \in I}\) be a family of sets, and for each \(i \in I\) let \(\mathcal{S}_i\) be a structure of species \(\Sigma\) on \(\mathbf{A}_i\) . Let E be the set which is the sum of the family \((\mathbf{A}_i)_{i \in I}\) , each \(\mathbf{A}_i\) being considered as a subset of E. Suppose that we are given a notion of \(\sigma\) -morphisms for the species \(\Sigma\) , and define an \(\alpha\) -mapping to be a mapping \(\varphi\) of E into a \(\Sigma\) -set F such that, for each \(i \in I\) , the restriction of \(\varphi\) to \(\mathbf{A}_i\) is a morphism of \(\mathbf{A}_i\) into F. Show that, if there exists a solution \((\mathbf{F}_E, \varphi_E)\) to the universal mapping problem for E, then the structure of species \(\Sigma\) on \(\mathbf{F}_E\) is the final structure with respect to the family \((\mathbf{A}_i, \mathcal{S}_i, \varphi_i)_{i \in I}\) , where \(\varphi_i\) denotes the restriction of \(\varphi_E\) to \(\mathbf{A}_i\) .
\end{enumerate}

Furthermore, let \(\mathbf{F}\) be a set, and for each \(\iota \in I\) let \(f_{\iota}\) be a mapping of \(\mathbf{A}_{\iota}\) into \(\mathbf{F}\) . If there exists a final structure of species \(\Sigma\) on \(\mathbf{F}\) with respect to the family \((\mathbf{A}_{\iota}, \mathcal{S}_{\iota}, f_{\iota})_{\iota \in I}\) , then we may write \(f_{\iota} = f \circ \varphi_{\iota}\) , where \(f\) is a morphism of \(\mathbf{F}_{\mathbf{E}}\) into \(\mathbf{F}\) , and the structure on \(\mathbf{F}\) is the direct image under \(f\) of the structure on \(\mathbf{F}_{\mathbf{E}}\) .

Consider the applications to the following cases :

* (i) \(\Sigma\) is a species of algebraic structures, the morphisms being homomorphisms, and the conditions \((\mathrm{CU}_{\mathbf{I}})\) through \((\mathrm{CU}_{\mathbf{III}})\) are satisfied. This is so for semi-group structures, group structures, module structures, algebra structures, etc. The \(\Sigma\) -set \(\mathbf{F}_{\mathbf{E}}\) is called the free product of the \(\mathbf{A}_{\mathbf{i}}\) in the case of groups, direct sum in the case of modules, direct product in the case of algebras.

\begin{enumerate}
\def\labelenumi{(\roman{enumi})}
\setcounter{enumi}{1}
\tightlist
\item
  \(\Sigma\) is the species of topological group structures or the species of topological vector space structures. The conditions \((\mathrm{CU}_{\mathrm{I}})\) through \((\mathrm{CU}_{\mathrm{III}})\) are then satisfied. In the case of locally convex topological vector spaces, \(\mathbf{F}_{\mathbf{E}}\) is called the topological direct sum of the \(\mathbf{A}_{i}\) . *
\end{enumerate}
